% ============================================================
%  CHAPITRE 4 — IMPLÉMENTATION TECHNIQUE
% ============================================================
\chapter{Implémentation technique}

\section{Introduction}

Ce chapitre présente l'implémentation concrète du système de maintenance prédictive.
Nous détaillons les technologies utilisées, l'implémentation de chaque composant,
l'intégration des modèles ML, les fonctionnalités principales, et les résultats
d'évaluation quantitatifs.

L'implémentation a suivi une approche itérative : chaque service a été développé et
testé indépendamment avant d'être intégré dans le système global. Cette approche a
permis de détecter et corriger les problèmes d'intégration au fur et à mesure du
développement.

% ─────────────────────────────────────────────────────────────
\section{Technologies utilisées}
% ─────────────────────────────────────────────────────────────

\begin{table}[H]
\centering
\caption{Stack technique complète du projet}
\begin{tabularx}{\textwidth}{|l|l|X|}
\hline
\rowcolor{lightblue}
\textbf{Couche} & \textbf{Technologie} & \textbf{Rôle} \\
\hline
Frontend & React 18.3 + Vite 5.4 & Interface utilisateur SPA \\
\hline
Frontend & TailwindCSS 3.4 & Design dark glassmorphism \\
\hline
Frontend & Recharts & Graphiques de métriques temporelles \\
\hline
Frontend & Axios & Appels API REST avec intercepteurs JWT \\
\hline
Backend & Node.js 20 + Express 4.21 & API REST centrale \\
\hline
Backend & Sequelize ORM & Accès PostgreSQL sécurisé (protection injection SQL) \\
\hline
Backend & JWT + bcrypt & Authentification et hashage des mots de passe \\
\hline
Backend & Nodemailer & Envoi d'emails d'alerte HTML via SMTP \\
\hline
Base de données & PostgreSQL 14 & Stockage persistant, verrous consultatifs \\
\hline
Agent & Python 3.9 + psutil & Collecte métriques système \\
\hline
Agent & smartctl & Lecture données SMART (HDD/SSD/NVMe) \\
\hline
Agent & schedule & Planification horaire autonome \\
\hline
ML Service & scikit-learn 1.7 & Random Forest + Isolation Forest \\
\hline
ML Service & PyTorch & Réseau LSTM (entraîné sur Backblaze) \\
\hline
ML Service & Flask + APScheduler & API ML + Planification cron (2h00) \\
\hline
ML Service & pandas + numpy & Traitement des séries temporelles \\
\hline
Chatbot & Ollama + LLaMA 3.2:1b & LLM local (0 cloud, 0 API externe) \\
\hline
Évaluation & rouge-score & Métriques ROUGE-1, ROUGE-2 (F1) \\
\hline
Déploiement & Docker + docker-compose & Conteneurisation 5 services \\
\hline
\end{tabularx}
\end{table}

% ─────────────────────────────────────────────────────────────
\section{Implémentation de l'agent de collecte}
% ─────────────────────────────────────────────────────────────

\subsection{Architecture modulaire de l'agent}

L'agent est composé de quatre modules Python avec des responsabilités clairement séparées.
Cette modularité permet de tester chaque composant indépendamment et de remplacer un
module sans affecter les autres.

\subsection{Collecte des métriques système — \texttt{SystemCollector}}

La méthode \texttt{collect\_all()} de la classe \texttt{SystemCollector} collecte les
trois catégories de métriques de manière résiliente : si une catégorie échoue (par exemple
les capteurs de température indisponibles), les autres continuent d'être collectées.

\textbf{CPU :} \texttt{psutil.cpu\_percent(interval=1)} effectue une mesure sur une
seconde pour éviter les valeurs instantanées non représentatives. La température est
recherchée séquentiellement dans les capteurs connus : \texttt{coretemp} (Intel),
\texttt{cpu\_thermal} (ARM), \texttt{k10temp} et \texttt{zenpower} (AMD). Si aucun
capteur n'est disponible (machines virtuelles, conteneurs), la valeur de fallback 50°C
est utilisée.

\textbf{RAM :} \texttt{psutil.virtual\_memory()} retourne directement le pourcentage
d'utilisation et la mémoire disponible, convertie en mégaoctets par division par
$1024^2$.

\textbf{Disque :} \texttt{psutil.disk\_usage()} est appelé sur \texttt{C:\textbackslash}
(Windows) ou \texttt{/} (Linux), retournant l'usage en pourcentage et l'espace libre
en mégaoctets.

\subsection{Lecture SMART — \texttt{SmartReader}}

Le module \texttt{SmartReader} implémente une stratégie de lecture robuste à plusieurs
niveaux. Sur Windows, il localise \texttt{smartctl} dans le PATH système ou dans les
chemins d'installation par défaut. Il tente ensuite plusieurs combinaisons de
périphériques (\texttt{/dev/sda} avec \texttt{-d nvme}, \texttt{/dev/sda} sans option,
\texttt{/dev/pd0}) jusqu'à obtenir une réponse valide.

L'analyse de la sortie de \texttt{smartctl -H} détermine le statut de santé global :
\texttt{PASSED} $\rightarrow$ GOOD, \texttt{FAILED} $\rightarrow$ CRITICAL. L'analyse
de \texttt{smartctl -A} extrait la température (recherche du mot-clé \texttt{temperature}
dans les lignes) et les compteurs d'erreurs (\texttt{raw\_read\_error\_rate},
\texttt{reallocated\_sector}).

\subsection{Mécanisme de retry avec backoff exponentiel}

La classe \texttt{DataSender} implémente un mécanisme de retry robuste pour gérer les
défaillances réseau transitoires :

\begin{table}[H]
\centering
\caption{Comportement du mécanisme de retry de l'agent}
\begin{tabular}{|c|c|l|}
\hline
\rowcolor{lightblue}
\textbf{Tentative} & \textbf{Délai} & \textbf{Condition} \\
\hline
1 & 0s & Première tentative \\
\hline
2 & 1s & Erreur HTTP 5xx ou timeout \\
\hline
3 & 2s & Erreur HTTP 5xx ou timeout \\
\hline
— & Abandon & Erreur HTTP 4xx ou 3 tentatives épuisées \\
\hline
\end{tabular}
\end{table}

Les erreurs 4xx ne déclenchent pas de retry car elles indiquent un problème de
configuration (token invalide, payload malformé) qui ne se résoudra pas avec une
nouvelle tentative. Le code définit \texttt{max\_retries = 3} avec des délais de 1s et 2s.

% ─────────────────────────────────────────────────────────────
\section{Implémentation du backend Node.js}
% ─────────────────────────────────────────────────────────────

\subsection{Structure MVC du projet}

Le backend Node.js est organisé selon le pattern MVC (Modèle-Vue-Contrôleur) :

\begin{itemize}
  \item \texttt{src/controllers/} : logique métier (auth, data, alerts, chatbot, ML, users) ;
  \item \texttt{src/routes/} : définition des endpoints REST et application des middlewares ;
  \item \texttt{src/models/} : modèles Sequelize (Machine, Agent, SystemMetrics, SmartData,
        User, Alert) ;
  \item \texttt{src/middleware/} : authentification JWT, validation des payloads, gestion
        d'erreurs centralisée, limitation de taille des requêtes ;
  \item \texttt{src/services/} : services transversaux (emailService, chatbotService,
        ollamaService).
\end{itemize}

\subsection{Endpoints API principaux}

\begin{table}[H]
\centering
\caption{Endpoints principaux de l'API REST}
\begin{tabularx}{\textwidth}{|l|l|X|}
\hline
\rowcolor{lightblue}
\textbf{Méthode} & \textbf{Route} & \textbf{Description} \\
\hline
POST & /api/auth/login & Authentification, retourne JWT \\
\hline
POST & /api/auth/register & Création de compte utilisateur \\
\hline
POST & /api/data & Réception métriques agent [Bearer] \\
\hline
GET & /api/machines & Liste machines + dernières prédictions [JWT] \\
\hline
GET & /api/machines/:id/metrics & Historique métriques paginé [JWT] \\
\hline
GET & /api/dashboard/overview & KPIs globaux (machines, alertes, risque moyen) [JWT] \\
\hline
GET & /api/alerts & Alertes actives triées par sévérité [JWT] \\
\hline
PUT & /api/alerts/:id/acknowledge & Accusé de réception [JWT] \\
\hline
PUT & /api/alerts/:id/resolve & Résolution alerte [JWT] \\
\hline
GET & /api/ml/predictions/:id & Prédiction RF machine [JWT] \\
\hline
GET & /api/ml/lstm/predict/:id & Prédiction LSTM machine [JWT] \\
\hline
POST & /api/ml/train & Déclencher entraînement RF [JWT] \\
\hline
POST & /api/chatbot & Requête chatbot [JWT] \\
\hline
GET & /api/users & Liste utilisateurs [JWT + Admin] \\
\hline
DELETE & /api/users/:id & Supprimer utilisateur [JWT + Admin] \\
\hline
\end{tabularx}
\end{table}

\subsection{Middleware d'authentification}

Le middleware \texttt{authenticateToken} extrait le token JWT de l'en-tête
\texttt{Authorization: Bearer <token>}, le vérifie avec \texttt{jwt.verify()} et
injecte le payload décodé dans \texttt{req.user}. Le middleware \texttt{requireAdmin}
vérifie ensuite que \texttt{req.user.role === 'admin'} pour les routes protégées.

\subsection{Service d'envoi d'emails}

Le service \texttt{emailService} utilise Nodemailer pour envoyer des emails HTML
formatés lors de la création d'alertes HIGH ou CRITICAL. L'email inclut le nom de
la machine, le niveau de risque, les probabilités à 7/14/30 jours, et les 5 features
les plus contributives à la prédiction.

% ─────────────────────────────────────────────────────────────
\section{Implémentation du service ML}
% ─────────────────────────────────────────────────────────────

\subsection{Pipeline d'extraction de caractéristiques}

La classe \texttt{FeatureExtractor} extrait \textbf{63 caractéristiques} par machine
à partir des 30 derniers jours de données stockées dans PostgreSQL. Les données sont
chargées via deux requêtes SQL paramétrées (protection contre l'injection SQL) :
une sur \texttt{system\_metrics} et une sur \texttt{smart\_data}, filtrées par
\texttt{machine\_id} et plage de dates.

\begin{table}[H]
\centering
\caption{Détail des 63 caractéristiques extraites par \texttt{FeatureExtractor}}
\begin{tabularx}{\textwidth}{|c|l|X|}
\hline
\rowcolor{lightblue}
\textbf{Nb} & \textbf{Type} & \textbf{Description} \\
\hline
45 & Statistiques glissantes & mean, median, std, min, max $\times$ 3 métriques $\times$ 3 fenêtres (24h, 168h, 720h) \\
\hline
3  & Tendances & Pente de régression linéaire sur 7 jours (\texttt{numpy.polyfit}) \\
\hline
3  & Taux de variation & Différence heure/heure pour CPU, RAM, disque \\
\hline
3  & Volatilité & Coefficient de variation (std/mean) \\
\hline
6  & SMART & temperature, read\_errors, write\_errors, health one-hot (3 valeurs) \\
\hline
3  & Temporel & hour\_of\_day, day\_of\_week, is\_weekend \\
\hline
\end{tabularx}
\end{table}

\subsection{Pipeline d'entraînement — \texttt{TrainingPipeline}}

La classe \texttt{TrainingPipeline} orchestre l'ensemble du processus d'entraînement
en cinq étapes :

\begin{enumerate}
  \item \textbf{Extraction des données} : \texttt{extract\_data(days=90)} charge les
        données des 90 derniers jours pour toutes les machines actives ;
  \item \textbf{Préparation} : \texttt{prepare\_training\_data()} génère les labels
        synthétiques (CPU $> 80\%$, RAM $> 85\%$, disque $> 90\%$ sur 30 jours $\rightarrow$
        label = 1) et effectue un split 80/20 avec stratification si possible ;
  \item \textbf{Entraînement} : \texttt{train\_model()} entraîne le Random Forest avec
        les hyperparamètres définis, avec un mécanisme de fallback vers des paramètres
        plus simples en cas d'échec ;
  \item \textbf{Sauvegarde} : \texttt{save\_and\_register\_model()} persiste le modèle
        sérialisé (pickle) dans le registre \texttt{ml\_models} ;
  \item \textbf{Activation conditionnelle} : \texttt{\_check\_and\_activate\_model()}
        compare l'accuracy du nouveau modèle avec le modèle actif. Le nouveau modèle
        n'est activé que si l'amélioration est $\geq 5\%$, évitant les régressions.
\end{enumerate}

\subsection{Modèle Random Forest — Hyperparamètres et justification}

\begin{table}[H]
\centering
\caption{Hyperparamètres du Random Forest et justification}
\begin{tabularx}{\textwidth}{|l|c|X|}
\hline
\rowcolor{lightblue}
\textbf{Paramètre} & \textbf{Valeur} & \textbf{Justification} \\
\hline
\texttt{n\_estimators} & 100 & Compromis variance/temps de calcul \\
\hline
\texttt{max\_depth} & 10 & Évite l'overfitting sur 63 features \\
\hline
\texttt{min\_samples\_split} & 5 & Régularisation des nœuds internes \\
\hline
\texttt{min\_samples\_leaf} & 2 & Évite les feuilles trop spécifiques \\
\hline
\texttt{class\_weight} & balanced & Compense le déséquilibre des classes (pannes rares) \\
\hline
\texttt{random\_state} & 42 & Reproductibilité des résultats \\
\hline
\texttt{n\_jobs} & -1 & Parallélisation sur tous les cœurs CPU \\
\hline
\end{tabularx}
\end{table}

\subsection{Labels synthétiques et leurs limites}

En l'absence de données réelles de pannes dans notre parc de 20 machines, les labels
d'entraînement sont générés par des règles à seuils métier :

\begin{itemize}
  \item CPU moyen sur 30 jours $> 80\%$ $\rightarrow$ label = 1 (panne probable) ;
  \item RAM moyenne sur 30 jours $> 85\%$ $\rightarrow$ label = 1 ;
  \item Disque moyen sur 30 jours $> 90\%$ $\rightarrow$ label = 1.
\end{itemize}

Cette approche constitue une limitation reconnue du système : les labels synthétiques
ne capturent pas la complexité réelle des défaillances matérielles. Le modèle LSTM
entraîné sur le dataset Backblaze (pannes réelles documentées) compense partiellement
cette limitation pour les données SMART.

\subsection{Inférence Random Forest — \texttt{Predictor}}

La classe \texttt{Predictor} implémente l'inférence en production. Pour chaque machine,
elle :
\begin{enumerate}
  \item Charge le modèle actif depuis le registre (avec cache en mémoire pour éviter
        les rechargements répétés) ;
  \item Extrait les features des 30 derniers jours via \texttt{FeatureExtractor} ;
  \item Appelle \texttt{model.predict\_proba(X)[0][1]} pour obtenir la probabilité
        de la classe positive (panne), valeur entre 0 et 1, convertie en pourcentage
        (0--100) avant stockage dans la table \texttt{predictions} ;
  \item Dérive les probabilités à 7j et 14j par facteurs d'échelle (0,70 et 0,85) ;
  \item Calcule le niveau de risque (LOW/MEDIUM/HIGH/CRITICAL) ;
  \item Extrait les 5 features les plus importantes via \texttt{feature\_importances\_}.
\end{enumerate}

\subsection{Modèle LSTM — Architecture PyTorch}

La classe \texttt{LSTMModel} est définie en PyTorch avec l'architecture suivante :

\begin{itemize}
  \item \texttt{nn.LSTM(input\_size=4, hidden\_size=32, batch\_first=True)} ;
  \item \texttt{nn.Linear(32, 1)} suivi d'une activation sigmoïde ;
  \item Inférence sur séquences de longueur \texttt{SEQ\_LEN = 5}.
\end{itemize}

Le modèle est chargé en mode singleton (\texttt{\_load\_model()}) pour éviter les
rechargements répétés lors des appels API. Il est évalué en mode \texttt{eval()} avec
\texttt{torch.no\_grad()} pour désactiver le calcul des gradients et optimiser la
vitesse d'inférence.

\subsection{Classification du risque LSTM}

La sortie sigmoïde du LSTM est convertie en niveau de risque selon les seuils suivants :

\begin{table}[H]
\centering
\caption{Seuils de classification du risque LSTM}
\begin{tabular}{|c|c|l|}
\hline
\rowcolor{lightblue}
\textbf{Probabilité} & \textbf{Niveau} & \textbf{Explication générée} \\
\hline
$> 0{,}70$ & HIGH & Risque élevé dû à des erreurs disque croissantes \\
\hline
$0{,}50 - 0{,}70$ & MEDIUM & Signes précoces de dégradation détectés \\
\hline
$\leq 0{,}50$ & LOW & Métriques dans la plage normale \\
\hline
\end{tabular}
\end{table}

Le modèle LSTM retourne trois niveaux de risque (HIGH, MEDIUM, LOW), contrairement au
Random Forest qui en retourne quatre (LOW, MEDIUM, HIGH, CRITICAL). La sortie sigmoïde
est une valeur brute entre 0 et 1, affichée directement dans la colonne
``Santé Disque (LSTM)'' du tableau de bord sans conversion en pourcentage.

\subsection{Planificateur de prédictions — APScheduler}

Le scheduler APScheduler utilise un \texttt{CronTrigger} pour déclencher le pipeline
de prédictions chaque nuit à 2h00. Le pipeline complet est :

\begin{enumerate}
  \item Acquisition du verrou distribué PostgreSQL (\texttt{pg\_try\_advisory\_lock}) ;
  \item Récupération de la liste des machines depuis la DB ;
  \item Pour chaque machine : extraction des features $\rightarrow$ inférence RF
        $\rightarrow$ stockage de la prédiction dans \texttt{predictions} ;
  \item Création d'alertes et envoi d'emails pour les risques HIGH/CRITICAL ;
  \item Libération du verrou (\texttt{pg\_advisory\_unlock}).
\end{enumerate}

En cas d'échec, un retry est planifié 1 heure plus tard. Un avertissement est émis si
le job dépasse 30 minutes d'exécution.

% ─────────────────────────────────────────────────────────────
\section{Implémentation du chatbot}
% ─────────────────────────────────────────────────────────────

\subsection{Architecture hybride à trois niveaux}

Le chatbot adopte une architecture hybride qui optimise la vitesse et la cohérence des
réponses. Selon le type de question détecté, trois chemins de traitement sont possibles :

\begin{figure}[H]
\centering
\includegraphics[width=0.95\textwidth]{chatbot_architecture.png}
\caption{Architecture hybride du chatbot à trois niveaux}
\end{figure}

\subsection{Détection d'intention par règles regex}

La détection d'intention est basée sur des règles regex appliquées à la question
normalisée en minuscules. Cette approche déterministe garantit des temps de réponse
inférieurs à 10ms pour les intentions reconnues, sans appel au LLM :

\begin{table}[H]
\centering
\caption{Règles de détection d'intention du chatbot}
\begin{tabularx}{\textwidth}{|l|X|l|}
\hline
\rowcolor{lightblue}
\textbf{Intention} & \textbf{Déclencheurs (mots-clés)} & \textbf{Action} \\
\hline
greeting & bonjour, salut, hello, bonsoir & Ollama (ton amical) \\
\hline
knowledge & mots-clés base de connaissances & Réponse directe (< 10ms) \\
\hline
alerts & alerte, critique, alert & Requête SQL sur \texttt{alerts} \\
\hline
high\_risk\_list & quelles/liste + risque/élevé & Requête SQL sur \texttt{predictions} \\
\hline
machine\_count & combien, nombre, total & COUNT(*) sur \texttt{machines} \\
\hline
machine\_status & PC-XXX ou nom machine & Recherche ILIKE dans \texttt{machines} \\
\hline
general & tout le reste & Ollama + contexte DB injecté \\
\hline
\end{tabularx}
\end{table}

\subsection{Base de connaissances}

La base de connaissances contient 16 entrées couvrant les thèmes suivants :
maintenance prédictive, probabilité de panne, alertes CRITICAL/HIGH, collecte des
métriques, données SMART, modèle ML, CPU à 90\%, mise à jour des prédictions,
réduction du risque, disque à 95\%, alertes email, tableau de bord, anomalies,
graphiques de santé, et les scores ROUGE du chatbot lui-même
(ROUGE-1 = 0,4548, ROUGE-2 = 0,3043).

\subsection{Intégration Ollama}

Pour les questions générales, le service \texttt{ollamaService} envoie une requête
POST à \texttt{http://localhost:11434/api/generate} avec un prompt enrichi du contexte
de la base de données (nombre de machines, alertes actives, risque moyen). Les
paramètres d'inférence (\texttt{temperature = 0.3}, \texttt{num\_predict = 150},
\texttt{top\_k = 20}, \texttt{top\_p = 0.85}, \texttt{repeat\_penalty = 1.1})
sont configurés pour produire des réponses concises et factuellement cohérentes.

% ─────────────────────────────────────────────────────────────
\section{Implémentation du frontend React}
% ─────────────────────────────────────────────────────────────

\subsection{Structure de l'application SPA}

L'application React est une SPA (Single Page Application) avec React Router v6.
La navigation est protégée par un composant \texttt{DashboardWrapper} qui vérifie
la présence d'un token JWT valide avant d'afficher les routes protégées.

\begin{table}[H]
\centering
\caption{Routes de l'application frontend}
\begin{tabularx}{\textwidth}{|l|l|X|}
\hline
\rowcolor{lightblue}
\textbf{Route} & \textbf{Composant} & \textbf{Accès} \\
\hline
/ & LandingPage & Public \\
\hline
/login & Login & Public \\
\hline
/signup & Signup & Public \\
\hline
/dashboard & Dashboard & JWT requis \\
\hline
/dashboard (modal) & MachineDetails & JWT requis \\
\hline
/models & ModelPerformance & JWT requis \\
\hline
/users & UserManagement & JWT + Admin \\
\hline
\end{tabularx}
\end{table}

\subsection{Composants principaux}

\begin{description}
  \item[\texttt{KPICards}] : affiche les indicateurs clés (machines totales, alertes
        actives, risque moyen) avec des couleurs codées par niveau de risque ;
  \item[\texttt{MachineList}] : tableau des 20 machines affichant le hostname, l'adresse IP,
        les métriques CPU/RAM/Disque, le score de santé, la \textbf{probabilité de panne RF
        à 30 jours} (valeur 0--100\% stockée directement en base), et la \textbf{santé disque
        LSTM} issue du modèle PyTorch entraîné sur Backblaze. La colonne ``Risque Panne''
        a été supprimée au profit de ces deux colonnes de prédiction complémentaires ;
  \item[\texttt{MachineDetails}] : détail d'une machine avec probabilités 7/14/30j,
        top 5 features contributives (depuis \texttt{contributing\_factors} JSON),
        et graphiques Recharts de l'évolution des métriques ;
  \item[\texttt{AlertsList}] : composant intégré dans le dashboard affichant les 6 alertes
        les plus récentes en lecture seule, avec code couleur par sévérité (CRITICAL/WARNING)
        et horodatage relatif ;
  \item[\texttt{Chatbot}] : interface conversationnelle flottante avec affichage des
        scores ROUGE et historique de la conversation ;
  \item[\texttt{ModelPerformance}] : métriques des modèles ML (accuracy, F1, versioning)
        avec historique des 7 versions sauvegardées.
\end{description}

\subsection{Export CSV}

Le module \texttt{exportCSV.js} permet d'exporter les données de métriques d'une machine
au format CSV. Il génère un fichier avec les colonnes timestamp, cpu\_usage,
memory\_usage, disk\_usage, et les données SMART disponibles, téléchargeable directement
depuis le navigateur via l'API \texttt{Blob}.

% ─────────────────────────────────────────────────────────────
\section{Résultats d'évaluation}
% ─────────────────────────────────────────────────────────────

\subsection{Évaluation du modèle LSTM sur le dataset Backblaze}

Le modèle LSTM a été entraîné et évalué sur le dataset Backblaze Q1 2020 :

\begin{table}[H]
\centering
\caption{Résultats d'évaluation du modèle LSTM (Backblaze Q1 2020)}
\begin{tabular}{|l|c|l|}
\hline
\rowcolor{lightblue}
\textbf{Métrique} & \textbf{Score} & \textbf{Interprétation} \\
\hline
Accuracy & 99,2\% & 99,2\% des prédictions sont correctes \\
\hline
ROC AUC & 1,000 & Séparation parfaite normal/panne \\
\hline
Précision (Panne) & 92\% & 92\% des pannes prédites sont réelles \\
\hline
Rappel (Panne) & 100\% & Aucune panne réelle manquée \\
\hline
F1-Score (Panne) & 96\% & Excellent équilibre précision/rappel \\
\hline
MAE & 0,0070 & Erreur absolue moyenne très faible \\
\hline
RMSE & 0,0769 & Erreur quadratique moyenne faible \\
\hline
\end{tabular}
\end{table}

Le \textbf{rappel de 100\%} est la métrique la plus importante dans ce contexte :
le modèle ne manque aucune panne réelle dans l'ensemble de test. Cela est directement
lié à l'utilisation de \texttt{pos\_weight = 50} dans la fonction de perte
\texttt{BCEWithLogitsLoss}, qui pénalise fortement les faux négatifs. Le ROC AUC de
1,000 indique une séparation parfaite entre les disques sains et les disques défaillants
sur le dataset Backblaze.

\subsection{Analyse de la matrice de confusion LSTM}

Sur l'ensemble de test du dataset Backblaze, la matrice de confusion est la suivante :

\begin{table}[H]
\centering
\caption{Matrice de confusion du modèle LSTM sur Backblaze}
\begin{tabular}{|l|c|c|}
\hline
\rowcolor{lightblue}
 & \textbf{Prédit : Normal} & \textbf{Prédit : Panne} \\
\hline
\textbf{Réel : Normal} & TN (très élevé) & FP (8\% des prédictions positives) \\
\hline
\textbf{Réel : Panne} & FN = 0 & TP (100\% des pannes détectées) \\
\hline
\end{tabular}
\end{table}

L'absence de faux négatifs (FN = 0) confirme que le modèle est calibré pour la
maintenance prédictive : il préfère générer quelques fausses alarmes (FP) plutôt que
de manquer une vraie panne.

\subsection{Importance des features — Random Forest}

L'analyse des \texttt{feature\_importances\_} du modèle Random Forest révèle que les
features les plus prédictives sont les statistiques glissantes sur la fenêtre de 720h
(30 jours), notamment :

\begin{enumerate}
  \item \texttt{disk\_usage\_mean\_720h} : usage moyen du disque sur 30 jours ;
  \item \texttt{cpu\_usage\_mean\_720h} : usage moyen du CPU sur 30 jours ;
  \item \texttt{memory\_usage\_mean\_720h} : usage moyen de la RAM sur 30 jours ;
  \item \texttt{disk\_usage\_trend\_168h} : tendance du disque sur 7 jours (pente) ;
  \item \texttt{cpu\_usage\_std\_720h} : volatilité du CPU sur 30 jours.
\end{enumerate}

Ces résultats confirment que les tendances à long terme (30 jours) sont plus prédictives
que les valeurs instantanées, ce qui justifie le choix de la fenêtre de 720h comme
fenêtre principale d'extraction de features.

\subsection{Évaluation du chatbot — Scores ROUGE}

L'évaluation a été conduite sur un dataset de 20 questions-réponses de référence en
français, couvrant les thèmes : machines à risque, alertes, modèles ML, données SMART,
et administration. Le pipeline \texttt{evaluate.py} utilise \texttt{rouge-score} avec
\texttt{use\_stemmer=True}.

\begin{table}[H]
\centering
\caption{Scores ROUGE détaillés par question (20 questions de référence)}
\begin{tabularx}{\textwidth}{|c|X|c|c|}
\hline
\rowcolor{lightblue}
\textbf{\#} & \textbf{Question} & \textbf{ROUGE-1} & \textbf{ROUGE-2} \\
\hline
1  & Quelles machines sont à risque élevé?          & 0,390 & 0,256 \\
\hline
2  & Combien de machines sont surveillées?           & 0,286 & 0,182 \\
\hline
3  & Montre-moi les alertes critiques               & 0,400 & 0,242 \\
\hline
4  & Quel est le niveau de risque de la machine?    & 0,625 & 0,522 \\
\hline
5  & Qu'est-ce que la maintenance prédictive?       & 0,375 & 0,261 \\
\hline
6  & Comment interpréter une probabilité de 80\%?  & 0,583 & 0,478 \\
\hline
7  & Que signifie une alerte CRITICAL?              & 0,429 & 0,350 \\
\hline
8  & Différence entre HIGH et CRITICAL?             & 0,389 & 0,118 \\
\hline
9  & Comment fonctionne la collecte?                & 0,622 & 0,326 \\
\hline
10 & Qu'est-ce que les données SMART?               & 0,400 & 0,316 \\
\hline
11 & Comment le modèle ML prédit les pannes?        & 0,391 & 0,227 \\
\hline
12 & Que faire si le CPU dépasse 90\%?              & 0,392 & 0,286 \\
\hline
13 & Quand les prédictions sont mises à jour?       & 0,350 & 0,158 \\
\hline
14 & Comment réduire le risque de panne?            & 0,390 & 0,205 \\
\hline
15 & Que signifie un disque à 95\%?                 & 0,458 & 0,261 \\
\hline
16 & Comment sont envoyées les alertes email?       & 0,468 & 0,222 \\
\hline
17 & Quelle est la fréquence de collecte?           & 0,524 & 0,300 \\
\hline
18 & Comment accéder au tableau de bord?            & 0,511 & 0,444 \\
\hline
19 & Qu'est-ce qu'une anomalie système?             & 0,512 & 0,390 \\
\hline
20 & Comment interpréter les graphiques?            & 0,600 & 0,542 \\
\hline
\rowcolor{lightblue}
\multicolumn{2}{|c|}{\textbf{MOYENNE}} & \textbf{0,4548} & \textbf{0,3043} \\
\hline
\end{tabularx}
\end{table}

\textbf{Interprétation :} Un score ROUGE-1 de 0,4548 indique que 45\% des mots des
réponses de référence apparaissent dans les réponses du chatbot. Les meilleures
performances sont obtenues sur les questions relatives au niveau de risque d'une machine
(ROUGE-1 : 0,625), à la collecte des métriques (ROUGE-1 : 0,622) et aux graphiques de
santé (ROUGE-1 : 0,600). Ces questions correspondent aux intentions traitées par requêtes
SQL déterministes, qui produisent des réponses plus proches des références.

Les scores plus faibles (questions 2, 8, 13) correspondent aux questions traitées par
le LLM Ollama, dont le vocabulaire opérationnel diffère des références conceptuelles.
ROUGE mesure le chevauchement lexical, pas la correction sémantique : les réponses du
chatbot sont factuellement correctes mais utilisent des formulations différentes.

\subsection{Métriques de performance système}

\begin{table}[H]
\centering
\caption{Métriques de performance du système en production}
\begin{tabular}{|l|c|}
\hline
\rowcolor{lightblue}
\textbf{Indicateur} & \textbf{Valeur mesurée} \\
\hline
Temps de réponse API moyen & $<$ 100ms \\
\hline
Temps de requête DB moyen & $<$ 50ms \\
\hline
Chargement frontend initial & $<$ 2s \\
\hline
Durée cycle de collecte agent & $\sim$ 3s \\
\hline
Machines surveillées & 20 \\
\hline
Enregistrements \texttt{system\_metrics} & 7,8 millions \\
\hline
Versions modèles RF sauvegardées & 7 \\
\hline
Consommation RAM de l'agent & $<$ 50 Mo \\
\hline
Usage CPU de l'agent (collecte) & $<$ 1\% \\
\hline
\end{tabular}
\end{table}

% ─────────────────────────────────────────────────────────────
\section{Interface utilisateur — Présentation des écrans}
% ─────────────────────────────────────────────────────────────

Cette section présente les principales vues de l'interface React développée dans le cadre
du projet. Chaque écran correspond à un besoin fonctionnel identifié lors de la phase de
conception et implémenté dans un composant React dédié.

% ── 1.1 Vue globale ──────────────────────────────────────────
\subsection{Vue globale du système (Dashboard)}

\begin{figure}[H]
  \centering
  \includegraphics[width=0.95\textwidth]{screen_dashboard.png}
  \caption{Vue globale du système — Tableau de bord principal}
\end{figure}

Le tableau de bord principal constitue le point d'entrée de l'application après
authentification. Il est rendu par le composant \texttt{Dashboard.jsx} et agrège les
données retournées par l'endpoint \texttt{GET /api/dashboard/overview}.

La vue affiche trois catégories d'informations :
\begin{itemize}
  \item \textbf{KPIs globaux} (composant \texttt{KPICards}) : nombre total de machines
        surveillées (20), nombre d'alertes actives, et niveau de risque moyen du parc
        exprimé en pourcentage. Les cartes sont colorées selon le niveau de risque
        (vert/jaune/orange/rouge) ;
  \item \textbf{Liste des machines} (composant \texttt{MachineList}) : tableau des 20
        machines avec leur hostname, adresse IP, métriques CPU/RAM/Disque, score de santé,
        \textbf{probabilité de panne Random Forest à 30 jours} (colonne ``Proba. Panne (RF)'',
        valeur stockée en 0--100\% dans la table \texttt{predictions}), et
        \textbf{santé disque LSTM} (colonne ``Santé Disque (LSTM)'', sortie sigmoïde du
        modèle PyTorch entraîné sur le dataset Backblaze). La colonne ``Risque Panne''
        a été retirée pour laisser place à ces deux indicateurs ML distincts ;
  \item \textbf{Alertes actives} (composant \texttt{AlertsList}) : liste des alertes
        HIGH et CRITICAL en attente d'accusé de réception.
\end{itemize}

% ── 1.1b Tableau des machines ────────────────────────────────
\subsection{Tableau des machines — Double prédiction ML}

\begin{figure}[H]
  \centering
  \includegraphics[width=0.95\textwidth]{screen_machinelist.png}
  \caption{Tableau \texttt{MachineList} — Colonnes ``Proba. Panne (RF)'' et ``Santé Disque (LSTM)''}
\end{figure}

Le composant \texttt{MachineList} affiche les deux prédictions ML côte à côte pour
chaque machine :
\begin{itemize}
  \item \textbf{Proba. Panne (RF)} : probabilité de panne à 30 jours calculée par le
        Random Forest, stockée en 0--100\% dans la table \texttt{predictions}. La couleur
        varie du vert ($< 30\%$) à l'orange ($\geq 50\%$) et au rouge ($\geq 70\%$) ;
  \item \textbf{Santé Disque (LSTM)} : sortie sigmoïde du modèle PyTorch entraîné sur
        le dataset Backblaze, valeur entre 0 et 1. Affiche ``N/A'' si moins de 5 mesures
        SMART sont disponibles.
\end{itemize}

\textbf{Analyse de cas réel — Machine ``Mori'' (PC de développement) :}
La machine réelle utilisée pour le développement du projet affiche une probabilité RF
de 86,8\% (CRITICAL) et une santé disque LSTM de 6\% (LOW). Ces deux valeurs, apparemment
contradictoires, illustrent parfaitement la complémentarité des deux modèles :

\begin{itemize}
  \item Le \textbf{Random Forest} détecte une saturation des ressources système (CPU,
        RAM, disque) due à l'exécution simultanée de tous les services de la plateforme
        (backend, ML service, frontend, PostgreSQL, Ollama) sur la même machine. Les
        labels synthétiques d'entraînement associent une utilisation élevée à un risque
        de panne, ce qui explique le score CRITICAL ;
  \item Le \textbf{LSTM Backblaze} analyse les attributs SMART du disque physique
        (secteurs réalloués, erreurs de lecture/écriture, température). Le disque étant
        matériellement sain, le modèle retourne un risque faible de 6\%.
\end{itemize}

Cette divergence démontre que les deux modèles mesurent des dimensions différentes du
risque : le RF évalue la \textit{charge opérationnelle}, tandis que le LSTM évalue la
\textit{dégradation matérielle}. Un technicien averti interprétera le score RF élevé
comme un signal de surcharge à corriger (répartition des services), et non comme une
défaillance imminente du matériel confirmée par le LSTM.

% ── 1.2 Analyse des machines ─────────────────────────────────
\subsection{Analyse détaillée d'une machine}

\begin{figure}[H]
  \centering
  \includegraphics[width=0.95\textwidth]{screen_machine_details.png}
  \caption{Analyse détaillée d'une machine — Prédictions et métriques}
\end{figure}

La vue de détail d'une machine est rendue par le composant \texttt{MachineDetails.jsx},
accessible en cliquant sur une machine dans le tableau de bord (modal). Elle présente :

\begin{itemize}
  \item \textbf{KPIs instantanés} : CPU, RAM et Disque en pourcentage avec code couleur ;
  \item \textbf{Prédiction LSTM} : jauge circulaire affichant la probabilité de panne
        (sortie sigmoïde du modèle PyTorch), le niveau de risque (LOW/MEDIUM/HIGH) et
        l'explication textuelle générée par le modèle. Un badge ``Anomalie'' s'affiche
        si \texttt{lstm.anomaly = true} ;
  \item \textbf{Anomalies détectées} : liste des anomalies Isolation Forest des 7
        derniers jours, avec type, métrique, valeur, plage attendue et score d'anomalie ;
  \item \textbf{Graphique Recharts} : évolution temporelle CPU/RAM/Disque sur 6h, 12h,
        24h ou 48h, sélectionnable via des boutons ;
  \item \textbf{Informations machine} : numéro de série, OS, statut, santé disque SMART ;
  \item \textbf{Actions admin} : boutons ``Modifier'' (hostname, IP) et ``Supprimer''
        (admin uniquement, avec confirmation).
\end{itemize}

% ── 1.3 Alertes & Prédictions ───────────────────────────────
\subsection{Alertes et prédictions (Dashboard)}

\begin{figure}[H]
  \centering
  \includegraphics[width=0.95\textwidth]{screen_alerts.png}
  \caption{Alertes et prédictions — Composant \texttt{AlertsList.jsx}}
\end{figure}

Le composant \texttt{AlertsList.jsx} est intégré directement dans le tableau de bord
principal. Il affiche un flux des alertes récentes générées par le pipeline ML,
triées par date de création (les plus récentes en premier).

Pour chaque alerte, les informations suivantes sont affichées :
\begin{itemize}
  \item Type d'alerte (CPU, MEMORY, DISK, SMART, PREDICTION) avec badge coloré ;
  \item Nom de la machine concernée et son adresse IP ;
  \item Message descriptif de l'alerte ;
  \item Horodatage relatif (ex. : ``Il y a 5 min'').
\end{itemize}

Le composant est en lecture seule — il affiche les 6 alertes les plus récentes
directement dans le dashboard sans nécessiter de navigation vers une page dédiée.

% ── 1.3b Email d'alerte ──────────────────────────────────────
\subsection{Notification par email (Alerte HIGH/CRITICAL)}

\begin{figure}[H]
  \centering
  \includegraphics[width=0.95\textwidth]{screen_email_alert.png}
  \caption{Email d'alerte automatique reçu pour une machine à risque HIGH}
\end{figure}

Lorsqu'une machine atteint un niveau de risque HIGH ($\geq 50\%$) ou CRITICAL
($\geq 70\%$), le backend génère automatiquement un email HTML via Nodemailer/SMTP
(Gmail). L'email contient :
\begin{itemize}
  \item Le niveau de sévérité avec code couleur (orange pour HIGH, rouge pour CRITICAL) ;
  \item Le nom et l'adresse IP de la machine concernée ;
  \item Les probabilités de panne à 7, 14 et 30 jours ;
  \item Un lien direct vers le tableau de bord (\texttt{http://localhost:3001}).
\end{itemize}

% ── 1.4 Assistant intelligent ────────────────────────────────
\subsection{Assistant intelligent (Chatbot)}

\begin{figure}[H]
  \centering
  \includegraphics[width=0.95\textwidth]{screen_chatbot.png}
  \caption{Assistant intelligent — Interface conversationnelle LLaMA}
\end{figure}

Le chatbot est implémenté dans le composant \texttt{Chatbot.jsx} sous la forme d'une
fenêtre flottante accessible depuis toutes les pages de l'application. Il communique
avec le backend via \texttt{POST /api/chatbot}.

L'interface affiche :
\begin{itemize}
  \item L'historique de la conversation avec distinction visuelle entre les messages
        de l'utilisateur (alignés à droite) et les réponses du système (alignés à gauche) ;
  \item Les scores ROUGE-1 et ROUGE-2 du chatbot (0,4548 et 0,3043), affichés en bas
        de la fenêtre pour informer l'utilisateur de la qualité évaluée du système ;
  \item Un indicateur de chargement pendant le traitement de la requête par Ollama.
\end{itemize}

L'architecture hybride est transparente pour l'utilisateur : les questions opérationnelles
(machines à risque, alertes actives) reçoivent des réponses instantanées via requêtes SQL,
tandis que les questions générales sont traitées par LLaMA 3.2:1b en local.

% ── 1.5 Évaluation des modèles ───────────────────────────────
\subsection{Évaluation des modèles ML (LSTM + Random Forest)}

\begin{figure}[H]
  \centering
  \includegraphics[width=0.95\textwidth]{screen_model_lstm.png}
  \caption{Onglet LSTM --- Métriques MAE / MSE / RMSE et historique des versions}
\end{figure}

\begin{figure}[H]
  \centering
  \includegraphics[width=0.95\textwidth]{screen_model_rf.png}
  \caption{Onglet Random Forest --- Accuracy, Précision, Rappel, F1-Score}
\end{figure}

Le composant \texttt{ModelPerformance.jsx} est accessible via le bouton `Modèles ML''
dans l'en-tête du dashboard. Il présente deux onglets distincts :

\textbf{Onglet LSTM (thème violet) :}
\begin{itemize}
  \item Métriques de régression : MAE = 0,0070, MSE = 0,0059, RMSE = 0,0769 ;
  \item Graphique à barres Recharts des trois métriques ;
  \item Historique des versions LSTM avec statut actif/inactif.
\end{itemize}

\textbf{Onglet Random Forest (thème vert) :}
\begin{itemize}
  \item Métriques de classification : Accuracy, Précision, Rappel et F1-Score,
        lus depuis la table \texttt{ml\_models} ;
  \item Graphique à barres avec axe Y en pourcentage (0--100\%) ;
  \item Historique des versions RF avec leurs métriques respectives.
\end{itemize}

Cette vue à deux onglets permet de comparer directement les deux modèles : le LSTM
évalué sur des métriques de régression (erreur de prédiction continue) et le Random
Forest évalué sur des métriques de classification (qualité de la décision binaire)
panne / pas de panne).
\subsection{Gestion des utilisateurs (Admin)}

\begin{figure}[H]
  \centering
  \includegraphics[width=0.95\textwidth]{screen_user_management.png}
  \caption{Gestion des utilisateurs — Interface administrateur}
\end{figure}

La vue de gestion des utilisateurs est rendue par le composant \texttt{UserManagement.jsx},
accessible uniquement aux utilisateurs avec le rôle \texttt{admin} via la route
\texttt{/users}. L'accès est protégé côté backend par le middleware \texttt{requireAdmin}.

Cette vue permet à l'administrateur de :
\begin{itemize}
  \item Consulter la liste complète des comptes utilisateurs avec leur email, rôle
        (admin/viewer) et date de création ;
  \item Créer de nouveaux comptes via un formulaire intégré (appel à
        \texttt{POST /api/auth/register}) ;
  \item Modifier le rôle d'un utilisateur existant ;
  \item Supprimer un compte utilisateur via \texttt{DELETE /api/users/:id}.
\end{itemize}

Les mots de passe ne sont jamais affichés ni transmis en clair : ils sont hashés avec
bcrypt (10 rounds) avant stockage dans la table \texttt{users}.

% ── 1.0 Landing Page ─────────────────────────────────────────
\subsection{Page d'accueil (Landing Page)}

\begin{figure}[H]
  \centering
  \includegraphics[width=0.95\textwidth]{screen_landing.png}
  \caption{Page d'accueil — \texttt{LandingPage.jsx}}
\end{figure}

La page d'accueil est le point d'entrée public de l'application, accessible via la route
\texttt{/} sans authentification. Rendue par le composant \texttt{LandingPage.jsx}, elle
présente le nom du système (\textbf{PC Technician Assistant}), une description synthétique
des fonctionnalités principales (surveillance en temps réel, prédiction ML, alertes
automatiques, chatbot IA), et deux boutons de navigation vers la connexion
(\texttt{/login}) et l'inscription (\texttt{/signup}).

% ── 1.7 Authentification ─────────────────────────────────────
\subsection{Authentification}

Le système d'authentification comprend trois écrans publics accessibles sans token JWT :
la connexion, l'inscription, et la réinitialisation du mot de passe.

\begin{figure}[H]
  \centering
  \includegraphics[width=0.95\textwidth]{screen_login.png}
  \caption{Page de connexion — \texttt{Login.jsx}}
\end{figure}

La page de connexion implémente le flux JWT complet :
\begin{enumerate}
  \item L'utilisateur saisit son email et son mot de passe ;
  \item Le frontend envoie \texttt{POST /api/auth/login} ;
  \item Le backend vérifie via \texttt{bcrypt.compare()}, génère un token JWT signé
        (\texttt{JWT\_SECRET}, expiration 24h) et le retourne ;
  \item Le frontend stocke le token et redirige vers \texttt{/dashboard}.
\end{enumerate}

\begin{figure}[H]
  \centering
  \includegraphics[width=0.95\textwidth]{screen_signup.png}
  \caption{Page d'inscription — \texttt{Signup.jsx}}
\end{figure}

La page d'inscription (\texttt{Signup.jsx}, route \texttt{/signup}) permet la création
de nouveaux comptes via \texttt{POST /api/auth/register}. Les champs sont validés côté
client (format email, longueur du mot de passe) et côté serveur par le middleware
\texttt{validation.js}. Le mot de passe est hashé avec bcrypt (10 rounds) avant stockage.

\begin{figure}[H]
  \centering
  \includegraphics[width=0.95\textwidth]{screen_forgot_password.png}
  \caption{Réinitialisation du mot de passe — \texttt{ForgotPassword.jsx}}
\end{figure}

La page de réinitialisation (\texttt{ForgotPassword.jsx}, route \texttt{/forgot-password})
permet à un utilisateur de demander la réinitialisation de son mot de passe via son
adresse email. Un email de réinitialisation est envoyé via le service Nodemailer.

% ─────────────────────────────────────────────────────────────
\section{Déploiement avec Docker}
% ─────────────────────────────────────────────────────────────

\subsection{Conteneurisation des services}

Le système est entièrement conteneurisé avec Docker et docker-compose. Chaque service
dispose de son propre Dockerfile optimisé pour la production :

\begin{table}[H]
\centering
\caption{Services Docker du système}
\begin{tabular}{|l|c|l|}
\hline
\rowcolor{lightblue}
\textbf{Service} & \textbf{Port exposé} & \textbf{Image de base} \\
\hline
postgres   & 5432 & postgres:14-alpine \\
\hline
backend    & 3000 & node:20-alpine \\
\hline
frontend   & 3001 & node:20-alpine \\
\hline
ml-service & 5000 & python:3.9-slim \\
\hline
agent      & —    & python:3.9-slim \\
\hline
\end{tabular}
\end{table}

Le déploiement complet se fait avec une seule commande : \texttt{docker-compose up -d}.
Les variables d'environnement sensibles (mots de passe, clés JWT, credentials SMTP)
sont injectées via des fichiers \texttt{.env} non versionnés.

\subsection{Gestion des dépendances entre services}

Le fichier \texttt{docker-compose.yml} définit les dépendances entre services via
\texttt{depends\_on} avec des conditions de santé (\texttt{healthcheck}) :
le backend attend que PostgreSQL soit prêt avant de démarrer, et le service ML attend
que le backend soit disponible. L'agent attend que le backend soit accessible avant
de commencer la collecte.

% ─────────────────────────────────────────────────────────────
\section{Conclusion}
% ─────────────────────────────────────────────────────────────

Ce chapitre a présenté l'implémentation complète du système, depuis l'agent de collecte
jusqu'au chatbot IA, en passant par les modèles ML et le tableau de bord. Les résultats
d'évaluation confirment la pertinence des choix techniques : le modèle LSTM atteint un
rappel de 100\% sur les pannes réelles Backblaze (ROC AUC = 1,000), et le chatbot obtient
un score ROUGE-1 de 0,4548 sur 20 questions de référence. Le système gère 7,8 millions
d'enregistrements avec des temps de réponse inférieurs à 100ms, démontrant la scalabilité
de l'architecture microservices adoptée.
